% ***************************************************************************************************
%
%	Szablon pracy magisterskiej dla Politechniki Wrocławskiej w wersji dwustronnej.
%	Autor:	Tomasz Strzałka
%
%	Udostępniony dzięki uprzejmości dr Macieja Gębali z Wydziału Podstawowych Problemów Techniki
%
% ***************************************************************************************************

% Styl dwustronny z domyślną wielkością czcionki 10pt oraz oddzieloną stroną tytułową (titlepage).
% Domyślnie rodziały rozpoczynają się na stronie prawej (openright).
\documentclass[12pt]{book}

% ***************************************************************************************************
% Ustawienia języka
% ***************************************************************************************************

\usepackage[english]{babel}

% Podstawowe ustawienia języka, według którego formatowany będzie dokument

% Pakiet babel dla polskiego języka powoduje konflikt z pakietem amssymb.
% Polecenie '\lll' definiują oba pakiety - porządana jest druga definicja.
\let\lll\undefined

% W przypadku wielojęzykowości ustawia główny język dokumentu
\selectlanguage{english}

% Kodowanie dokumentu
\usepackage[utf8]{inputenc}

% Dowolny rozmiar czcionek, kodowanie znaków
\usepackage{lmodern}

% Polskie wcięcia akapitów
\usepackage{indentfirst}

% Polskie łamanie wyrazów
% \usepackage[plmath]{polski}

% Przecinek w wyrażeniach matematycznych zamiast kropki
% \usepackage{textualicomma}
\usepackage[T1]{fontenc}

% Polskie formatowanie typograficzne
\frenchspacing

% Zapewnia liczne usprawnienia wyświetlania i organizacji matematycznych formuł. 
\usepackage{amsmath}

% Wprowadza rozszerzony zestaw symboli m.in. \leadsto
\usepackage{amssymb}

% Dodatkowa, ,,kręcona'' czcionka matematyczna
\usepackage{mathrsfs}

% Dodatkowe wsparcie dla środowiska mathbb, które nie wspiera domyślnie cyfr (\mathbb{})
\usepackage{bbold}

% Fixes/improves amsmath
\usepackage{mathtools}

% adds floating barier
\usepackage[section]{placeins}

% long tables
\usepackage{longtable}

% backslash in table
\usepackage{diagbox}

% uncomment to show line numbers
\usepackage{lineno}
% \linenumbers

% Typesets in-line quotations
\usepackage{dirtytalk}

% ***************************************************************************************************
% Kolory  
% ***************************************************************************************************

% Umożliwia kolorowanie poszczególnych komórek tabeli
\usepackage[table]{xcolor}% http://ctan.org/pkg/

% Umożliwia łatwą zmianę koloru linii w tabeli
\usepackage{tabu}

% Umożliwia rozszerzoną kontrolę nad kolorami.
\usepackage{xcolor}

% Definicje kolorów
\definecolor{lgray}{HTML}{9F9F9F}
\definecolor{dgray}{HTML}{5F5F5F}
\definecolor{wit-red}{HTML}{9A1E13}
% lgray				-	nazwa nowo zdefiniowanego koloru
% HTML				-	model kolorów
% CCCCCC			-	wartość koloru zgodna z modelem

% ***************************************************************************************************
% Algorytmy 
% ***************************************************************************************************

% Udostępnia środowisko do konstruowania pseudokodów
\usepackage[ruled,vlined,linesnumbered,longend,algochapter]{algorithm2e}
% ruled	- poziome kreski na początku i końcu algorytmu, podpis na górze oddzielony również kreską poziomą
% vlined - pionowe kreski łączące początek polecenia z jego końcem
% linesnumbered	- numerowanie kolejnych wierszy algorytmu
% longend - długie końcówki np. ifend, forend itd.
% algochapter - numeracja z rozdziałami

% Zamiana nazwy środowiska z domyślnej "Algorithm X" na "Pseudokod X"
% \newenvironment{pseudokod}[1][htb]{
% 	\renewcommand{\algorithmcfname}{
% Algorithm}
% 	\begin{algorithm}[#1]%
% 	}{
% \end{algorithm}
% }

% Zmiana rozmiaru komentarzy
\newcommand\algcomment[1]{
	\footnotesize{#1}
}

% Ustawienie zadanego stylu dla komentarzy
\SetCommentSty{algcomment}

% Wyśrodkowana tylda
\usepackage{textcomp}%
\newcommand{\textapprox}{\raisebox{0.5ex}{\texttildelow}}

% Listowanie kodów źródłowych
\usepackage{listings} 
\renewcommand{\lstlistingname}{Source Code} % Polska nazwa listingu

% ***************************************************************************************************
% Inne 
% ***************************************************************************************************

% Umieszczanie notatek typu "todo"
\usepackage{todonotes} 

% Obsługa przypadków w równaniach matematycznych
\usepackage{cases}

% Pozwala na użycie komendy `\include` bez automatycznego wstawiania nowej strony
\usepackage{newclude}

% Pozwala na automatyczne zliczanie sekcji
\usepackage{totcount}
\regtotcounter{chapter}

\renewcommand{\lstlistingname}{Source Code} % Polska nazwa listingu
% Definicje pecjalnych znaków, które nie są obsługiwane w środowisku listing
\lstset{literate=
	{ż}{{\.{z}}}1	{ź}{{\'{z}}}1
	{ć}{{\'{c}}}1	{ń}{{\'{n}}}1
	{ą}{{\c a}}1	{ś}{{\'{s}}}1
	{ł}{{\l}}1		{ę}{{\c{e}}}1
	{ó}{{\'{o}}}1	{á}{{\'a}}1
	{é}{{\'e}}1		{í}{{\'i}}1
	{ó}{{\'o}}1		{ú}{{\'u}}1
	{ù}{{\`u}}1		{Á}{{\'A}}1
	{É}{{\'E}}1		{Í}{{\'I}}1
	{Ó}{{\'O}}1		{Ú}{{\'U}}1
	{à}{{\`a}}1		{è}{{\'e}}1
	{ì}{{\`i}}1		{ò}{{\`o}}1
	{ò}{{\`o}}1		{À}{{\`A}}1
	{È}{{\'E}}1		{Ì}{{\`I}}1
	{Ò}{{\`O}}1		{Ò}{{\`O}}1
	{ä}{{\"a}}1		{ë}{{\"e}}1
	{ï}{{\"i}}1		{ö}{{\"o}}1
	{ü}{{\"u}}1		{Ä}{{\"A}}1
	{Ë}{{\"E}}1		{Ï}{{\"I}}1
	{Ö}{{\"O}}1		{Ü}{{\"U}}1
	{â}{{\^a}}1		{ê}{{\^e}}1
	{î}{{\^i}}1		{ô}{{\^o}}1
	{û}{{\^u}}1		{Â}{{\^A}}1
	{Ê}{{\^E}}1		{Î}{{\^I}}1
	{Ô}{{\^O}}1		{Û}{{\^U}}1
	{œ}{{\oe}}1		{Œ}{{\OE}}1
	{æ}{{\ae}}1		{Æ}{{\AE}}1
	{ß}{{\ss}}1		{ç}{{\c c}}1
	{Ç}{{\c C}}1	{ø}{{\o}}1
	{å}{{\r a}}1	{Å}{{\r A}}1
	{€}{{\EUR}}1	{£}{{\pounds}}1
}

% ***************************************************************************************************
% Marginesy 
% ***************************************************************************************************

% Pozbycie się pustych stron po rozdziałach przy pomocy redfinicji.
\let\cleardoublepage\clearpage

% Ustawienia rozmiarów stron i ich marginesów
\usepackage[headheight=18pt, top=25mm, bottom=25mm, left=25mm, right=25mm]{geometry}
% headheight		-	wysokość tytułów
% top				-	margines górny
% bottom			-	margines dolny
% left				-	margines lewy
% right				-	margines prawy

% Usunięcie górnego marginesu dla środowisk
\makeatletter
\setlength\@fptop{0\p@}	
\makeatother

% ***************************************************************************************************
% Styl 
% ***************************************************************************************************

% Definiuje środowisko 'titlingpage', które zapewnia pełną kontrolę nad układem strony tytułowej.
\usepackage{titling}


% Umożliwia modyfikowanie stylu spisu treści
\usepackage{tocloft}	

\tocloftpagestyle{tableOfContentStyle}

% Definiowanie własnych stylów nagłówków i/lub stopek
\usepackage{fancyhdr}

% Domyślny styl dla pracy 
\fancypagestyle{custom}{
	\fancyhf{}									% wyczyść stopki i nagłówki
	\fancyhead[RO]{								% Prawy, nieparzysty nagłówek
		\hrulefill \hspace{16pt} \large Chapter \thechapter
		\put(-472.1, 12.1){%
			\makebox(0,0)[l]{%
				\includegraphics[width=0.05\textwidth]{images/pwr-logo.png}
			}
		}
		\put(-443,5.5){%
			\makebox(0,0)[l]{%
				\small Wrocław University of Science and Technology
			}
		}
	}
	\fancyhead[LE]{								% Lewy, parzysty nagłówek
		\large Chapter \thechapter \hspace{16pt} \hrulefill   
		\put(-17.5, 10){%
			\makebox(0,0)[l]{%
				\includegraphics[width=0.03\textwidth]{images/wit-logo.png}
			}
		}
		\put(-287.5, 8.5){%
			\makebox(0,0)[l]{%
				\small Faculty of Information and Communication Technology
			}
		}
	}
	\fancyfoot[LE,RO]{							% Stopki
		\thepage
	}
	\renewcommand{\headrulewidth}{0pt}			% Grubość linii w nagłówku
	\renewcommand{\footrulewidth}{0.2pt}		% Grubość linii w stopce
}


% Domyślny styl dla bibliografii
\fancypagestyle{bibliographyStyle}{
	\fancyhf{}									% wyczyść stopki i nagłówki
	\fancyhead[RO]{								% Prawy, nieparzysty nagłówek
		\hrulefill \hspace{16pt} \large Bibliography \thechapter
		\put(-472.1, 12.1){%
			\makebox(0,0)[l]{%
				\includegraphics[width=0.05\textwidth]{images/pwr-logo.png}
			}
		}
		\put(-443,5.5){%
			\makebox(0,0)[l]{%
				\small Wrocław University of Science and Technology
			}
		}
	}
	\fancyhead[LE]{								% Lewy, parzysty nagłówek
		\large Bibliography \hspace{16pt} \hrulefill  
		\put(-17.5, 10){%
			\makebox(0,0)[l]{%
				\includegraphics[width=0.03\textwidth]{images/wit-logo.png}
			}
		}
		\put(-287.5, 8.5){%
			\makebox(0,0)[l]{%
				\small Faculty of Information and Communication Technology
			}
		}
	}
	\fancyfoot[LE,RO]{							% Stopki
		\thepage
	}
	\renewcommand{\headrulewidth}{0pt}			% Grubość linii w nagłówku
	\renewcommand{\footrulewidth}{0.2pt}		% Grubość linii w stopce
}

% Domyślny styl dla dodatków
\fancypagestyle{appendixStyle}{
	\fancyhf{}									% wyczyść stopki i nagłówki
	\fancyhead[RO]{								% Prawy, nieparzysty nagłówek
		\hrulefill \hspace{16pt} \large Appendix \thechapter
		\put(-472.1, 12.1){%
			\makebox(0,0)[l]{%
				\includegraphics[width=0.05\textwidth]{images/pwr-logo.png}
			}
		}
		\put(-443,5.5){%
			\makebox(0,0)[l]{%
				\small Wrocław University of Science and Technology
			}
		}
	}
	\fancyhead[LE]{								% Lewy, parzysty nagłówek
		\large Appendix \thechapter \hspace{16pt} \hrulefill   
		\put(-17.5, 10){%
			\makebox(0,0)[l]{%
				\includegraphics[width=0.03\textwidth]{images/wit-logo.png}
			}
		}
		\put(-287.5, 8.5){%
			\makebox(0,0)[l]{%
				\small Faculty of Information and Communication Technology
			}
		}
	}
	\fancyfoot[LE,RO]{							% Stopki
		\thepage
	}
	\renewcommand{\headrulewidth}{0pt}			% Grubość linii w nagłówku
	\renewcommand{\footrulewidth}{0.2pt}		% Grubość linii w stopce
}

% Osobny styl dla stron zaczynających rozdział/spis treści itd. (domyślnie formatowane jako "plain")
\fancypagestyle{chapterBeginStyle}{
	\fancyhf{}%
	\fancyfoot[LE,RO]{
		\thepage
	}
	\renewcommand{\headrulewidth}{0pt}
	\renewcommand{\footrulewidth}{0.2pt}
}

% Styl dla pozostałych stron spisu treści
\fancypagestyle{tableOfContentStyle}{
	\fancyhf{}%
	\fancyfoot[LE,RO]{
		\thepage
	}
	\renewcommand{\headrulewidth}{0pt}
	\renewcommand{\footrulewidth}{0.2pt}
}

% Formatowanie tytułów rozdziałów i/lub sekcji
\usepackage{titlesec}

% Formatowanie tytułów rozdziałów
\titleformat{\chapter}[hang]					% kształt
{
	\vspace{-15ex}
	\Huge
	\bfseries
}												% formatowanie tekstu modyfikowanego elementu
{
    \thechapter	
}												% etykieta występująca przed tekstem modyfikowanego elementu, niewidoczna w spisie treści
{
	10pt
}												% odstęp formatowanego tytułu od lewego marginesu/etykiety
{
	\Huge
	\bfseries
}												% formatowanie elementów przed modyfikowanym tytułem
[
\vspace{2ex}
%\rule{\textwidth}{0.4pt}
%\vspace{-4ex}
]												% dodatkowe formatowanie stosowane poniżej modyfikowanego tytułu


% Formatowanie tytułów sekcji
\titleformat{\section}[hang]					% kształt
{
	\vspace{2ex}
%	\titlerule\vspace{1ex}
	\LARGE\bfseries
}												% formatowanie tekstu modyfikowanego elementu
{
	\thesection									% etykieta występująca przed tekstem modyfikowanego elementu, niewidoczna w spisie treści
}
{
	0pt
}												% odstęp formatowanego tytułu od lewego marginesu/etykiety
{
	\LARGE
	\bfseries
}												% formatowanie elementów przed modyfikowanym tytułem

% Formatowanie tytułów podsekcji
\titleformat{\subsection}					% kształt
{
	\Large\bfseries
}												% formatowanie tekstu modyfikowanego elementu
{
	\thesubsection									% etykieta występująca przed tekstem modyfikowanego elementu, niewidoczna w spisie treści
}
{
	0pt
}												% odstęp formatowanego tytułu od lewego marginesu/etykiety
{
	\Large
	\bfseries
}												% formatowanie elementów przed modyfikowanym tytułem


% Formatowanie tytułów podpodsekcji
\titleformat{\subsubsection}[hang]					% kształt
{
	\large\bfseries
}												% formatowanie tekstu modyfikowanego elementu
{
	\thesubsubsection									% etykieta występująca przed tekstem modyfikowanego elementu, niewidoczna w spisie treści
}
{
	0pt
}												% odstęp formatowanego tytułu od lewego marginesu/etykiety
{
	\large
	\bfseries
}												% formatowanie elementów przed modyfikowanym tytułem
% Formatowanie tytułów podpodsekcji
\titleformat{\paragraph}[hang]					% kształt
{
	\normalfont\bfseries
}												% formatowanie tekstu modyfikowanego elementu
{
	\theparagraph									% etykieta występująca przed tekstem modyfikowanego elementu, niewidoczna w spisie treści
}
{
	0pt
}												% odstęp formatowanego tytułu od lewego marginesu/etykiety
{
	\normalfont
	\bfseries
}												% formatowanie elementów przed modyfikowanym tytułem

% ***************************************************************************************************
% Linki
% ***************************************************************************************************

% Umożliwia wstawianie hiperłączy do dokumentu
\usepackage{hyperref}							% Aktywuje linki

\hypersetup{
	colorlinks	=	true,					% Koloruje tekst zamiast tworzyć ramki.
	linkcolor	=	wit-red,					% Kolory: referencji,
    citecolor	=	wit-red,					% cytowań,
	urlcolor	=	wit-red					% hiperlinków.
}

% Do stworzenia hiperłączy zostanie użyta ta sama (same) czcionka co dla reszty dokumentu
\urlstyle{same}




% ***************************************************************************************************
% Linki
% ***************************************************************************************************

% Umożliwia zdefiniowanie własnego stylu wyliczeniowego
\usepackage{enumitem}

% Nowa lista numerowana z trzema poziomami
\newlist{myitemize}{itemize}{3}

% Definicja wyglądu znacznika pierwszego poziomu
\setlist[myitemize,1]{
	label		=	\textbullet,
	leftmargin	=	4mm}

% Definicja wyglądu znacznika drugiego poziomu
\setlist[myitemize,2]{
	label		=	$\diamond$,
	leftmargin	=	8mm}

% Definicja wyglądu znacznika trzeciego poziomu
\setlist[myitemize,3]{
	label		=	$\diamond$,
	leftmargin	=	12mm
}

% ***************************************************************************************************
% Inne pakiety
% ***************************************************************************************************

% Dołączanie rysunków
\usepackage{graphicx}
\usepackage{wrapfig}

% Wykresy
\usepackage{pgfplots}

% Figury i przypisy
\usepackage{caption}
\usepackage{subcaption}

% Umożliwia tworzenie przypisów wewnątrz środowisk
\usepackage{footnote}

% Umożliwia tworzenie struktur katalogów
\usepackage{dirtree}

% Rozciąganie komórek tabeli na wiele wierszy
\usepackage{multirow}

% Precyzyjne obliczenia szerokości/wysokości dowolnego fragmentu wygenerowanego przez LaTeX
\usepackage{calc}

% Umożliwia automatyczne przeliczanie i wyświetlanie liczb w systemie rzymskim
\usepackage{romannum}

% Obsługa obrazów w formacie .svg
\usepackage{svg}

% ***************************************************************************************************
% Matematyczne skróty
% ***************************************************************************************************

% Skrócony symbol liczb rzeczywistych
\newcommand{\RR}{\mathbb{R}}

% Skrócony symbol liczb naturalnych
\newcommand{\NN}{\mathbb{N}}

% Skrócony symbol liczb wymiernych
\newcommand{\QQ}{\mathbb{Q}}

% Skrócony symbol liczb całkowitych
\newcommand{\ZZ}{\mathbb{Z}}

% Skrócony symbol logicznej implikacji
\newcommand{\IMP}{\rightarrow}

% Skrócony symbol  logicznej równoważności
\newcommand{\IFF}{\leftrightarrow}


% ***************************************************************************************************
% Środowiska
% ***************************************************************************************************

% Środowisko do twierdzeń
\newtheorem{theorem}{Theorem}[chapter]

% Środowisko do lematów
\newtheorem{lemma}{Lemma}[chapter]

% Środowisko do przykładów
\newtheorem{example}{Example}[chapter]

% Środowisko do wniosków
\newtheorem{corollary}{Conclusion}[chapter]

% Środowisko do definicji
\newtheorem{definition}{Definition}[chapter]
\providecommand*{\definitionautorefname}{Definition}

% Środowisko do dowodów
\newenvironment{proof}{
	\par\noindent \textbf{Proof.}
}{
\begin{flushright}
	\vspace*{-6mm}\mbox{$\blacklozenge$}
\end{flushright}
}

%multicolumn
\usepackage{multicol}

% Środowisko do uwag
\newenvironment{remark}{
	\bigskip \par\noindent \small \textbf{Remark.}
}{
\begin{small}
	\vspace*{4mm}
\end{small}
}

% ***************************************************************************************************
% Słownik
% ***************************************************************************************************

% Prawidłowe dzielenie wyrazów
% \hyphenation{wszy-stkich ko-lu-mnę każ-da od-leg-łość
% 	dzie-dzi-ny dzie-dzi-na rów-nych rów-ny
% 	pole-ga zmie-nna pa-ra-met-rów wzo-rem po-cho-dzi
% 	o-trzy-ma wte-dy wa-run-ko-wych lo-gicz-nie
% 	skreś-la-na skreś-la-ną cał-ko-wi-tych wzo-rów po-rzą-dek po-rząd-kiem
% 	przy-kład pod-zbio-rów po-mię-dzy re-pre-zen-to-wa-ne
% 	rów-no-waż-ne bi-blio-te-kach wy-pro-wa-dza ma-te-ria-łów
% 	prze-ka-za-nym skoń-czo-nym moż-esz na-tu-ral-na cią-gu tab-li-cy
% 	prze-ka-za-nej od-po-wied-nio}

% ***************************************************************************************************
% Dokument
% ***************************************************************************************************


\usepackage[english, masters]{title_page_wiz}

% Set up the thesis
\author{Michał Karbownik}
\title{Representation Learning with Bayesian Methods}
\supervisor{dr hab. inż.Tomasz Kajdanowicz}
\fieldofstudy{Artificial Intelligence}
\specialty{}
\keywords{%TODO
Bayesian representation learning, Bayesian, representation, deep learning, uncertainty}
\summary{%TODO 
This work tries to bridge a gap between recent representation learning advancements and their lack of proper Bayesian perspective by introducing, analyzing, and testing a cross-section of Bayesian representation learning models by an appropriate theoretical introduction, along with an observational and quantitative study.}
\year=2022

% \frontmatter

\pgfplotsset{compat=1.17} 
\begin{document}

\maketitle

	\section*{Abstract}
Learning representations of the data that make it easier to extract useful information when building classifiers or other predictors is referred to as representation learning. Recent years have brought a tremendous advancement in this area of Artificial Intelligence. It seems, however, as if this development neglected the benefits Bayesian approach could provide, namely uncertainty quantification, interpretability, or better latent space properties resulting in more universal and quality representations. This work tries to bridge this gap by introducing, analyzing, and testing a cross-section of Bayesian representation learning models by an appropriate theoretical introduction, along with an observational and quantitative study.

\section*{Streszczenie}
Uczenie się reprezentacji danych, które ułatwiają wydobywanie przydatnych informacji podczas budowania klasyfikatorów lub innych predyktorów, jest nazywane uczeniem reprezentacji. Ostatnie lata przyniosły ogromny postęp w tej dziedzinie sztucznej inteligencji. Wydaje się jednak, że rozwój ten pomija korzyści, jakie może zapewnić podejście bayesowskie, a mianowicie kwantyfikację niepewności, interpretowalność lub lepsze właściwości przestrzeni ukrytej, co skutkuje bardziej uniwersalnymi i jakościowymi reprezentacjami. Niniejsza praca próbuje wypełnić tę lukę, przedstawiając, analizując i testując przekrój modeli uczenia reprezentacji bayesowskich poprzez odpowiednie wprowadzenie teoretyczne, wraz z badaniami obserwacyjnymi oraz ilościowymi.
	
        \cleardoublepage
		
	\pagenumbering{Roman}
	\pagestyle{tableOfContentStyle}
	\tableofcontents
	\cleardoublepage
		
	% ***************************************************************************************************
	% Wstęp
	% ***************************************************************************************************
	
	\pagestyle{custom}
	\mainmatter
	
    \listoffigures
    
    \listoftables
	
	% ***************************************************************************************************
	% Rodziały
	% ***************************************************************************************************

	\chapter{Introduction}\label{chapter:intro}
\thispagestyle{chapterBeginStyle}
The idea of representations, while coming from the neuroscience background, has found its adoption in the field of Artificial Intelligence in a form of \textit{representation learning}. The learned representations should be transferable to other tasks and demonstrate certain qualities, such as limited dimensionality, embedding hidden features required for knowledge transfer, being interpretable or benefiting other predictors in their downstream tasks.

\vspace{\baselineskip}
Especially recent years have brought advancement in this area. It seems, however, as if this development neglected the benefits Bayesian approach could provide. There is little to no work considering this perspective. This work introduces the problem of representation learning from a Bayesian angle.

\section{Motivation}
As already mentioned, literature on the topic of Bayesian representation learning is scarce. The existing work is usually task-specific (usually with publicly unavailable data), which prevents reproducibility and confines application opportunities. Additionally, available methods are often very novel and complex, keeping the entry threshold for utilizing them high. Moreover, no previous work has been found that explicitly defined this problem, gathered the methods and introduced the topic of Bayesian representation learning.

\section{Aim and scope}\label{sec:scope}
The main objective of this work is to introduce, analyze and test a cross-section of Bayesian representation learning methods. In order to meet the objective, the following scope was defined:

\begin{itemize}
    \item Provide a literature review and theoretical background of feature learning and Bayesian approach to it.
    \item Provide theoretical background for 3 Bayesian representation learning methods: Bayesian PCA, Bayesian Neural Network and Variational Autoencoder.
    \item Analyze those methods in an observational study, highlighting their characteristics.
    \item Perform downstream task testing on representations learned with those methods.
\end{itemize}

\section{Research problems}
In this work, the following problems were identified and investigated:
\begin{itemize}
    \item How do PCA, neural networks, autoencoders, and their Bayesian counterparts function, and how do they obtain representations?
    \item What are possible advantages and drawbacks of using Bayesian PCA, Bayesian Neural Networks and Variational Autoencoders for representation learning?
    \item What characteristics do those methods demonstrate?
    \item How well do representations learned by them perform on downstream tasks with other classifiers?
\end{itemize}

\section{Contribution}
While this work by no means is an exhaustive survey of existing methods, it aims to introduce, apply and test a cross-section of Bayesian representation learning methods. A by-product of this is an introduction to Bayesian representation learning, which, with adequate adjustments, could be used in self-study or didactics. By testing fundamental methods of this paradigm, it provides a background and lays the foundation for the future work.


\section{Work structure}
This work consists of \total{chapter} chapters. Below, a short summary of each of them is provided.

\begin{enumerate}
    \item \autoref{chapter:intro} is an introductory chapter. It explains the motivation and contribution of this work, along with enumerating its aim and stated research problems.
    \item \autoref{chapter:background} provides necessary theoretical introduction and background on representation learning and Bayesian approach to it.
    \item \autoref{chapter:methods} explains the theory and mechanics of Bayesian PCA, Bayesian Neural Networks and Variational Autoencoders, with a mention of their conventional counterparts. 
    \item \autoref{chapter:experiments} describes used datasets and experimental setup.
    \item \autoref{chapter:observational} conducts an observational study on the three Bayesian representation learning methods.
    \item \autoref{chapter:quantitative} performs a series of downstream task quantitative experiments on trained representations.
    \item \autoref{chapter:conclusion} draws conclusions and provides ideas for possible future work.
\end{enumerate}
% 	\cleardoublepage

	\chapter{Background and literature review}\label{chapter:background}
\thispagestyle{chapterBeginStyle}
\label{theoretical-background}

\section{Representation learning}
\textit{Representation learning}, often referred to as \textit{feature learning}, is a set of machine learning approaches that allows for automatic feature detection from data. This is in contrast to manual feature extraction which, while takes advantage of the domain knowledge about data and human ingenuity, is both time-consuming, labor-intensive and cannot take advantage of the latent space. A comprehensible definition of representation learning was formulated by Bengio et al. \cite{Bengio2012} in 2012:

    \begin{definition}[Representation learning]\label{def:representation-learning}
    Learning representations of the data that make it easier to extract useful information when building classifiers or other predictors.
    \end{definition}
\autoref{def:representation-learning} does not, however, delve into the core issue of what the representation is – even though the definition seems to be intuitive, it is hardly to be found in the domain-specific literature.

\vspace{\baselineskip}
\subsection{What is a representation?}\label{subsec:representation}
As it often comes with Artificial Intelligence, being highly inspired by processes taking place in nature (in particular by the field of cognition, which is the main interest of cognitive science), it is not the sole domain whose interests revolve around representations. Not only the term itself, but also the very ontology of representations, play a central role in neuroscience literature. However, even there neither a clear consensus on all aspects, nor a unified definition of it exist \cite{Vilarroya2017}. Having said that, probably the most established is \autoref{def:representation-neuroscience}.

\begin{definition}[Representation (in neuroscience)]\label{def:representation-neuroscience}
A formal system for making explicit certain entities or types of information, together with a specification of how the system does this. \cite{Marr1982}
\end{definition}

An example of such would be a numeral system, which is a writing system for expressing (representing) numbers — the Roman, Hindu–Arabic (decimal) and binary being among the most widespread. \autoref{fig:numeric_systems} shows different representations of the same number in these systems.

% TODO: Create captions for particular representations.
% TODO: Improve main caption.
\begin{figure}[ht]
    \centering
    \includegraphics[width=0.8\textwidth]{background/images/numeric_systems.png}
    \caption[Numeric systems as representations]{Numeric systems as representations.}
    \label{fig:numeric_systems}
\end{figure}

% TODO: Find reference for decimal & binary being positional
% TODO: Find reference for binary being foundation of modern computing
Each of those representations was formed with a specific goal in mind – to simplify certain operations. The Roman system, used by the ancient Romans to count and perform other day-to-day transactions. It is non-positional \cite{Hazewinkel1990} and governed by principles of subtraction and addition. Despite being quite primitive, it had quickly become overcomplicated due to its poor scalability \cite{Ifrah1998}. It constituted a barrier to the mathematical development of the Roman Empire that had been significant to the earlier Arabic cultures. They had used the decimal numeric system, which, being multiplicative by nature (easily decomposed into powers of $10$) and positional, facilitates more complex arithmetics. The binary system, using only two symbols - $0$ and $1$ - through advancements of the Boolean algebra applied to the electrical switches, is widely used in modern computers. In general, each of them provides a different kind of abstraction to aid a specific objective.

\vspace{\baselineskip}
Despite the information content being the same in all three cases, their representation is different, which implies the orthogonality of the two. Switching between representations should not influence the original value they describe, and the transition process should be unambiguous due to explicitness (see \autoref{def:representation-neuroscience}). 
Even though systems might share the same symbols, values expressed by them might differ (i.e., a $1$ on $n$-$th$ position will produce a different value in decimal and binary system). Because of that, a representation should be defined by the shape of the manifold on which the data lie within the representational space, as opposed to a single piece of information \cite{Thompson2018}.

\vspace{\baselineskip}
According to the above, almost anything can be subject to being represented, from everyday objects to complex abstract concepts. Similarly, representation learning touches upon different aspects and applications of machine learning (for particular examples, see \autoref{sec:related-work}). A useful data representation should yield an improvement of the performance of a model, as well as facilitate the learning process – the scale of both highly depends on appropriately selected features (it is especially true for the deep learning models).
% TODO: Find why it is especially true for the deep models (literature).

\vspace{\baselineskip}
With this perspective, one could form a loose definition of \textit{representation} in machine learning:

\begin{definition}[Representation]\label{def:representation-ml}
A function mapping input data to a latent vector, which is useful for a task.
\end{definition}

Subsequently, \textit{representation learning} can be comprehended as an optimization problem of parameters of such function. If $x$ denotes the data, $\theta$ the parameters and $\vec{z}$ - the latent vector, the whole can expressed as \autoref{eq:representation-learning}.

\begin{equation}\label{eq:representation-learning}
    f\left(x;\theta\right)\rightarrow\vec{z}
\end{equation}
We expect the dimension of the representation to be not higher than the original data ($dim\left(\vec{z}\right) \le dim\left(x\right)$), but ideally it should be much lower ($dim\left(\vec{z}\right) \ll dim\left(x\right)$).


\subsection{What makes a good representation?}
According to \autoref{def:representation-ml}, the output set of features (ergo the latent vector) should be \textit{useful} for a specific \textit{task}. The task objective serves as a falsifiability mechanism defining their usefulness. It allows for eliminating uninformative features, and thus, reducing the dimensionality. The examples in \autoref{subsec:representation} emphasize the importance of a \textit{quality} representation. Hence, it is valid to ask: \say{\textit{what} makes a good representation?}. Surprisingly, it is an ill-defined and challenging question.


\subsubsection{Dimensionality reduction}
Yet again, cognitive science lays foundations for artificial intelligence and aids in answering fundamental questions. In 1959, Horace Barlow, identified the significance of \textit{redundancy reduction}\footnote{Barlow revisited the topic of redundancy reduction in 2001 \cite{Barlow2001}. He stated that \say{the original work was wrong in over-emphasizing the role of compressive coding and economy in neuron numbers, but right in drawing attention to the importance of redundancy}. Both documents accentuate the impact of Claude Shannon's work, as well as the \say{entanglement} and the mutual influence of neuroscience and artificial intelligence.} in information processing by human brains \cite{Barlow1959}. The original idea is to \say{achieve economy without losing any information at all}. It can be inexactly (representation learning usually assumes some information loss) translated to a dimensionality constraint defined along \autoref{eq:representation-learning}. Due to a threat of the \textit{curse of dimensionality}, it is desirable that the representation has minimal possible number of dimensions, while maintaining the variance level (at least for the data that is informative for the task at hand). Barlow's findings later became a base for research on disentangled representations. 


\subsubsection{Disentanglement}
The human perception of the world around heavily depends on visual recognition. Representations on the early stage of visual processing are not beneficial for the object recognition, because they are - \say{like two sheets of paper crumpled together} - \textit{tangled} \cite{DiCarlo2007}. Only at later stages it is being untangled \cite{DiCarlo2012}, therefore transformed from \say{visual representations that are easy to build, but are not easily decoded, into representations that we do not yet know how to build, but are easily decoded}.

\vspace{\baselineskip}
Correspondingly, building a solid representation for machine learning algorithms expects separation of information about the attributes (some features vary with respect to only one attribute), simplifying setting decision boundaries for subsequent tasks. Additionally, one wants to discard the build up of an unnecessary invariance related to nuisance attributes, but without disposing of practical information about the data. It diminishes the representation's sensitivity to the data's directions of variance that are not informative to the task. This procedure, called \textit{disentangling factors of variation} \cite{Szabo2017}, is nontrivial, but yields a representation, which is more robust than the original, tangled structure. 


\subsubsection{Geometric priors}
% TODO:
Despite real-world data often being incomplete (or even damaged), high-dimensional and chaotic, it often shows certain, usually hidden, regularities. A good representation should identify and capture them, ideally in a low-dimensional form. This limits the exponential space of possible functions. Such regularities can be exposed through unified geometric principles called \textit{geometric priors} \cite{Bronstein2021}. Two fundamental principles applied to physically structured data are scale separation and symmetry.

\vspace{\baselineskip}
\textit{Symmetry} of an object or a system is a transformation that leaves some property of the object or system unchanged or maintains its invariance. \textit{Scale separation}, on the other hand, retains important features of a signal during the transfer onto a coarser domain. 

\vspace{\baselineskip}
Invariance to scale, translation, or rotation often proves to be the key to improving performance of machine learning models. It is ubiquitous (and easily visualizable) in image processing. Convolutional Neural Networks, for instance, do incorporate the idea of geometric priors. Translational symmetry is achieved through filters with shared weights, while pooling exploits scale separation. They are not, however, fully invariant to rotation, unless appropriately extended.

\subsubsection{Smoothness}
Smoothness of the representation assumes its invariance to minor, insignificant changes of the data. Formally, it is expressed that if some input $x$ is approximately equal to some other input $y$ ($x \approx y$), the functions of both should also be approximately equal ($f(x) \approx f(y)$). It is one of the general priors, not only for representation learning, but also for the majority of machine learning approaches.

\vspace{\baselineskip}
Bengio et al. \cite{Bengio2012} argue that, while useful, this assumption alone is insufficient to overcome the \textit{curse of dimensionality} and \say{advocate learning algorithms that are flexible and non-parametric but do not rely exclusively on the smoothness assumption}.


\subsubsection{Distribution}
Distribution of a representation assumes spread of the representation's particular features across multiple parts of the same model — following the example of neural networks, the final feature value consisting of contributions from different neural paths. Moreover, it also induces feature abstraction and re-use, which facilitate utilization of simple representation to create more abstract ones in hierarchical models.


\subsection{Representation learning taxonomy}
%TODO introduce

\subsubsection{shallow vs. deep}\label{subsubsec:conventional-deep}
\textit{Deep learning} models, through multi-layer processing, facilitate learning multiple levels of abstraction of data representations. As per the definition of godfathers of deep learning – Yann LeCun, Yoshua Bengio and Geoffrey Hinton – they are \say{representation-learning methods with multiple levels of representation, obtained by composing simple but non-linear modules that each transform the representation at one level (starting with the raw input) into a representation at a higher, slightly more abstract level} \cite{LeCun2015}. In other words, the design of deep learning implicitly assumes learning representations. As opposed to this, the term \textit{shallow representation learning} (also referred to as \say{conventional}\footnote{Please note that in this work, the term \textit{conventional} is often used as an opposition to \textit{Bayesian}.} or \say{traditional}) will include all other methods.

\subsubsection{linear vs. nonlinear}
Linear models are such that adopt an idea of a linear transformation, which can be expressed by the following equation:
\begin{equation}
\label{eq:linear-model}
    y = X\beta + \epsilon
\end{equation}
where $y$ is a vector of observed values of a \textit{dependent variable} (output of the model), $X$ - a matrix consisting of vectors of \textit{independent variables}, a parameter vector $\beta$ and an unobserved random variable $\epsilon$ responsible for handling noise.

\vspace{\baselineskip}
In representation learning, it translates to linear separability, ergo independence, of particular features — while a convenient assumption, it limits the models' performance on data with nonlinearities that are often encountered in practical applications. Nonlinear models, on the other hand, while more expensive computationally, are capable of capturing those dependencies to some extent. A flagship example of a nonlinear approach is deep learning, mentioned in \autoref{subsubsec:conventional-deep}, however, the literature offers numerous extensions to linear models that allow for capturing nonlinear patterns.


\subsubsection{global vs. local}\label{subsec:taxonomy-global-local}
Global feature learning methods, aim to retain global knowledge about data within the learned feature space. For instance, PCA uses orthogonal linear transformation to change the coordinate system in such a way that retains as much variance as possible. It results in a set of linearly uncorrelated variables (converted from a higher-dimensional set of conceivably correlated ones).

\vspace{\baselineskip}
Such an approach, however, does not retain the data's local similarities, unlike local methods that address the issue by trying to discover the data's hidden manifold structure (thus they are sometimes referred to as \textit{manifold learning}).

\subsubsection{supervised vs. unsupervised}
\textit{Supervised learning} has emerged as one of the first and the most common machine learning paradigms. The reason for it is that our brains use it as one of the fundamental approaches to discovering the world, as well as developing and maintaining a variety of functions \cite{Knudsen1994}. It is characterized by utilizing annotated, ideally very large datasets. The task is to teach models how to yield the desired output (labels). In representation learning, it is expected for such a model to generalize well, in order to be able to use the learned features in other downstream tasks or in different models.

\vspace{\baselineskip}
In \textit{unsupervised learning} (this term was also derived from cognitive science \cite{Barlow1989}), the dataset is provided without the annotated labels. Its goal is to discover hidden patterns of data and cluster it without imposed boundaries (for instance by labels, like in supervised learning). Unsupervised learning used in representation learning is supposed to reduce dimensionality, uncover exploratory factors, and train vectors without human intervention.

\vspace{\baselineskip}
It is also important to mention two other paradigms that emerged in recent years, namely \textit{self-supervised} and \textit{semi-supervised} representation learning. The latter can handle and make use of data that is labeled only partially. It performs well when it has a large amount of them and only a small portion that is labeled \cite{Engelen2020}. In self-supervised learning, models are provided with unstructured input data, and their task is to generate the labels automatically. Those are further used in subsequent iterations as ground truths. This approach is characterized by robustness to label corruption (and adversarial examples) \cite{Hendrycks2019}, as well as generalization ability and data efficiency \cite{Liu2020}. 

\subsubsection{generative vs. contrastive}
Another taxonomy, introduced by Liu et al. \cite{Liu2020}, assumes the existence of \textit{generative} and \textit{contrastive} methods. The former's task is to produce a latent vector, from which they reconstruct the input, thus being able to generate data. The latter bases on an idea of measuring similarity between samples and  requires negative sampling, which often demands heavy data augmentation.

\section{Bayesian approach}
% Bayesian approach 
% http://allendowney.github.io/ThinkBayes2/chap01.html
% http://jakevdp.github.io/blog/2014/03/11/frequentism-and-bayesianism-a-practical-intro/
% 9.2 tfof https://arxiv.org/pdf/2002.08791.pdf
% Methods that are not considered Bayesian may still define a probability distribution
% on the data. However, Bayesian methods can use the probability distribution to express uncertainty about the data and, in most cases, model parameters.
The term \textit{Bayesian} has been a buzzword in machine learning for quite some time. As with every trend, it has also been overused and its initial meaning distorted. While by no means is it a complete guide to Bayesian statistics, this section provides necessary background knowledge to understand Bayesian representation learning. 


\subsection{Bayesian vs frequentist}
There is an ongoing debate in the scientific community around statistics. Sometimes it is so tempestuous that one could call it a dispute, or even an argument. The disagreement revolves around the superiority of one of two approaches, and can be summarized in one question: \say{Bayesian or frequentist?}.

\vspace{\baselineskip}
The latter, is an interpretation, in which the probability\footnote{In this work the term \textit{probability distribution} is often shortened to both \textit{probability} or \textit{distribution}. While in statistics they represent different concepts (probability distribution is a mathematical function that describes the probability of different possible values of a variable), in artificial intelligence they are often used interchangeably.} is strictly related to the frequency of events. It assumes the existence of true, unknown parameters, whose values are fixed. In this case, the inference is made over hypothetical random samples of data. In the context of machine learning, it manifests itself as optimization (rather than marginalization) and usually results in a point estimate of parameters achieved by different methods, of which the most common is \textit{Maximum Likelihood Estimation} (MLE).

% \todo{(?) Add an illustration for distribution vs single point}

\vspace{\baselineskip}
Bayesian inference, on the other hand, treats the model parameters as random variables and makes probabilistic statements about their distribution. At the very foundations, it is driven by the \textit{Bayes Theorem}, which is often expressed by \autoref{eq:bayes-theorem}:

\begin{equation}\label{eq:bayes-theorem}
    P(H \mid D) = \frac{P(H)P(D \mid H)}{P(D)},
\end{equation}

\noindent where $H$ is called a \textit{hypothesis}, and $D$ is data.

\vspace{\baselineskip}
This formula can be easily adopted to be used in a model, like in \autoref{eq:bayes-theorem-model}:

\begin{subequations}
    \begin{equation}\label{eq:bayes-theorem-model}
        \overbrace{P(\theta \mid D, I)}^{\text{posterior}} = \frac{\overbrace{P(\theta \mid I)}^{\text{prior}}\overbrace{P(D \mid \theta, I)}^{\text{likelihood}}}{\underbrace{P(D \mid I)}_{\text{marginal likelihood}}}
    \end{equation}
    \begin{equation}\label{eq:bayes-theorem-model-marginalization}
        P(D \mid I) = \int P(D \mid \theta, I)P(\theta \mid I)d\theta
    \end{equation}
\end{subequations}

\noindent In this case, hypothesis ($H$) is replaced by parameters ($\theta$), and additional knowledge about the model (e.g., its type or hyperparameters) - $I$ is used. The \textit{prior distribution} specifies an assumption about the parameters without taking the data into account. The \textit{likelihood} represents the probability of the data under the specified parameters. The \textit{marginal likelihood} (also called \textit{evidence}) is the distribution of the data given our additional assumption. It can also be denoted without the parameters $\theta$ being marginalized out (\autoref{eq:bayes-theorem-model-marginalization}). Marginal distribution is the normalization of the Bayes rule and plays an important role in model comparison, however, is often intractable. Finally, the \textit{posterior distribution} is the distribution of the parameters after \say{seeing} the data (and the additional assumption $I$). 

\vspace{\baselineskip}
The above is sometimes simplified to the form: \say{the posterior is proportional to the likelihood and the prior} and denoted as:

\begin{equation}
    posterior \propto likelihood \times prior
\end{equation}
% Christopher Bishop in \textit{Pattern Recognition and Machine Learning} \cite{Bishop2006} states that \say{at the heart of Bayesian methods lie marginalizations}.
% The foundation of Bayes's Theorem, and thus the foundation of Bayesian statistics, is conditional probability.
% Conditional probablility -> Bayes Theorem -> equation -> inverse of thinking -> popularity

\subsection{Bayesian inference in representation learning}\label{subsec:bayesian-inference}
% https://www.youtube.com/watch?v=I4dkEALQv34
% https://www.youtube.com/watch?v=0w_4QcvBYII
In Bayesian representation learning, latent vectors $\mathbf{z}$ are not maximum a posteriori point estimations, but rather realizations from a posterior distribution, obtained by approximating the intractable marginal distribution.

% https://stats.stackexchange.com/questions/271844/variational-inference-versus-mcmc-when-to-choose-one-over-the-other

% https://towardsdatascience.com/bayesian-inference-problem-mcmc-and-variational-inference-25a8aa9bce29
In general, there are two types of approximation techniques — stochastic and deterministic. The most common examples of them in Bayesian setting are, respectively, sampling, especially \textit{Markov Chain Monte Carlo} (MCMC), and \textit{Variational Inference} (VI).


\paragraph{Markov Chain Monte Carlo}
The former, Markov Chain Monte Carlo, is a family of methods that, in theory, is able to provide exact sampling from the required distribution. In practice, however, it is not possible due to finite resources. The results achieved can be reliable when the sample is large enough. This is, in turn, computationally intensive; hence these methods are usually used for small-scale tasks with limited data. As the name suggests, it draws from two principles - \textit{Markov Chain} and \textit{Monte Carlo}.

\vspace{\baselineskip}
The latter is a concept of performing a random walk by drawing \textit{independent and identically distributed} (i.i.d) samples from a target distribution and estimating the desired quantity (such as variance or mean of the drawn samples). Introducing the \say{Markov Chain} part, however, weakens the i.i.d. constraint by requiring the dependency of the subsequent variables in a way that a given variable is dependent on its direct predecessor, which satisfies the Markov property:

\begin{equation}
    p(x^i \mid x^{i-1}, \dots, x^1) = p(x^i \mid x^{i-1}).
\end{equation}

Furthermore, MCMC methods require a phase called \textit{burn-in}, because they are sensitive to their starting point. During this phase, the initial samples are removed until the Markov chain converges to its stationary distribution, which adds to the computational intensity of the task.

\paragraph{Variational Inference}\label{paragraph:vi}
Variational inference methods are typically faster than MCMC because, instead of trying to approximate the target distribution, they redefine the problem to optimization with a simpler alternative (called \textit{variational distribution}). The variational distribution must be flexible and complex enough to provide a reasonable approximation to the true posterior, while being simple enough to facilitate optimization. Because, in fact, there is a difference between the distributions, it needs to be quantified and considered in the inference task. 

\vspace{\baselineskip}
\textit{Kullback-Leibler divergence} (KL or $D_{KL}$) is a score that measures how much information is lost when one distribution is approximated with another (and because of that, it is sometimes called \textit{relative entropy}). It fulfills the basic criteria for being such a measure:
% \begin{equation}
\begin{subnumcases}{}
    D_{KL}(p \parallel q) \geq 0\label{eq:kl-positive}\\ 
    D_{KL}(p \parallel q) = 0, & \textrm{if}\;p(x) = q(x)\\
    D_{KL}(p \parallel q) \neq KL(q \parallel p)
\end{subnumcases}
% \end{equation}

\noindent Kullback-Leibler divergence in the considered case can be defined as:
\begin{equation}\label{eq:kl-divergence}
    D_{KL}\left[q(\theta) \parallel p(\theta \mid D)\right] = \int q(\theta) \log\frac{q(\theta)}{p(\theta \mid D)}\, d\theta.
\end{equation}

\noindent The task is to minimize the KL divergence between the true and variational distributions, or more precisely, optimize the parameters of the latter. Because the unwanted term $p(\theta \mid D)$ is present in \autoref{eq:kl-divergence}, the Kullback-Leibler divergence is still intractable directly. A concept that can aid it, by transforming it into an optimization problem, is \textit{Evidence Lower Bound} (ELBO). Compared to the log-evidence of the likelihood function:

\begin{equation}\label{eq:elbo-iniequality}
    log\:p(D; \theta) \geq \mathbb{E}_q \left[ log \frac{p(D, \theta)}{q(\theta)} \right],
\end{equation}

% \begin{equation}\label{eq:kl-divergence-transformation}
% \begin{split}
%     \theta_{opt} &= \text{arg min}_\theta D_{KL}\left[q(\theta) \parallel p(\theta \mid D) \right] \\
%     &= \text{arg min}_\theta \int q(\theta) \log\frac{q(\theta)}{p(\theta)p(D \mid \theta)}\, d\theta\\
%     &= \text{arg min}_\theta D_{KL}\left[q(\theta) \parallel p(\theta) \right] - {\mathbb{E}}_{q(\theta)}\left[log\:p(D \mid \theta)\right],
% \end{split}
% \end{equation}
\noindent which, considering the aforementioned, can be transformed into:
\begin{equation}\label{eq:kl}
    D_{KL}\left[q(\theta) \parallel p(\theta \mid D) \right] = log\:p(D) - \mathbb{E}_q\left[log\:p(D, \theta) - log\:q(\theta) \right]
\end{equation}

\noindent where the term on the right-hand side is ELBO, and the other is evidence. It is important to note that because the evidence does not depend on $q$, it can be assumed constant. Furthermore, as in \autoref{eq:kl-positive}, the KL is always non-negative. Hence, a tractable operation of maximizing the ELBO is equivalent to minimizing the KL divergence, which allows for the inference.
% \subsection{Aleatoric vs Epistemic Uncertainty}
% https://arxiv.org/pdf/1703.04977.pdf

% https://alexgkendall.com/computer_vision/bayesian_deep_learning_for_safe_ai/

% https://towardsdatascience.com/building-a-bayesian-deep-learning-classifier-ece1845bc09



\vspace{\baselineskip}
\section{Bayesian representation learning}\label{sec:related-work}
If one were to determine the precursor of representation learning, it would undoubtedly be Principal Component Analysis. The groundwork for it was laid in 1901 by Karl Pearson \cite{Pearson1901}. Later, in 1936, Harold Hotelling \cite{Hotelling1936} named and developed it to its prevailing form. It was not until 1998 that its Bayesian counterpart was introduced by Christopher Bishop \cite{Bishop1998}. It is also this method that could be considered one of the first Bayesian representation learning methods.

\vspace{\baselineskip}
Representation learning in general was not widely popular until ca. 2006, when both the advancement of technology facilitated the development of deep learning. Only then the progress on feature learning techniques, which are often computationally expensive, could also be advanced.

\vspace{\baselineskip}
Even though theory for the Bayesian approaches has been available for a long time (a specific case of Bayes' theorem, which laid the foundation for the Bayesian statistics, was published already in 1763 \cite{Bayes1763}), it was not until a few decades ago that these methods were actually applied in machine learning \cite{Bishop1998, MacKay1992, Mackay1995} and they are being rediscovered only recently \cite{Jospin2022, Blundell2015, Wilson2020, Kingma2015, Chen2012, Gal2016, Kingma2014, Wang2018}.

\vspace{\baselineskip}
For the Bayesian representation learning, the field is even more scarce. Among the aforementioned Bayesian methods, there exist those, that enable representation learning, however, they are usually not studied from that perspective. In fact, only a few works discuss them, and even if — they are either task-specific (thus hardly universal) or very complex novel approaches \cite{Nozawa2020, Meng2021, Barkan2021, Zavatone2021}. 

% 	\cleardoublepage

	\chapter{Methods introduction}\label{chapter:methods}
\thispagestyle{chapterBeginStyle}
% \section{Representational similraity}
% https://www.researchgate.net/publication/306860121_Representational_models_A_common_framework_for_understanding_encoding_pattern-component_and_representational-similarity_analysis
% https://github.com/ZitongLu1996/NeuroRA
% https://github.com/rsagroup/rsatoolbox
% https://github.com/Charestlab/pyrsa
A quality theoretical introduction is a part of every quality dissertation. In order to use an apparatus, it is crucial to understand how and why it works. This chapter introduces the methods that will be used in the further study. It provides theoretical background, as well as necessary comments.

\include*{methods/bpca}
\include*{methods/bnn}
\include*{methods/vae}
% 	\cleardoublepage
	
	\chapter{Experimental setting}\label{chapter:experiments}
\thispagestyle{chapterBeginStyle}
Comparing different approaches does not always come down only to quantitative research. It is especially true for the ones providing more value than simply improving a score of a given metric. This is exactly the case of Bayesian methods — their main goal is usually not to compete with other models, but to provide additional information by capturing uncertainty. Due to the above, the conducted experiments focused on assessing not only the quantitative aspect, but also a comparative analysis in a broader sense.

\section{Datasets}\label{sec:datasets}
In order to minimize the number of variables during experiments, the chosen datasets were kept similar, both in terms of the appliable downstream tasks, but also attributes (such as sample size, modality, etc.). For the results to be reliable and not susceptible to bias, however, datasets, while being akin, need to provide enough variance. As proved by Cheplygina and Tax \cite{Cheplygina2018}, this is usually the case, even for datasets from the same source.

\vspace{\baselineskip}
Besides, it was desirable for the datasets to be well-established in the community. The reason is twofold: thorough knowledge about the dataset and avoiding unnecessary issues related to incorrectly created bespoke datasets (while there might be doubts about their flawlessness, such as annotation issues of MNIST, they are still much more robust and tested than experiment-specific ones).

\vspace{\baselineskip}
Another crucial characteristics was the size of the dataset. The ones with too few samples would not be enough to train and test representations reliably. The chosen dataset needed to be large enough to eliminate this risk.

%TODO: Zacytować prezentację z NIPS 2021


\vspace{\baselineskip}
The choice was to select only image classification datasets, since it brings even more advantages, which the specificity of conducted experiments required, such as simple and intuitive visualizations of the learning process' intermediate steps, or broader space of possible experiments (e.g. translation or rotation invariance).

\vspace{\baselineskip}
Three datasets with gradual complexity were chosen. This allowed for thorough examination of tested models, no matter their performance, and resulted in a more cross-sectional study. The following were used:



\vspace{\baselineskip}
%TODO Add visualisations of MNIST and CIFARs
\begin{itemize}
    \item MNIST \cite{deng2012mnist, LeCun1998}
    \\\\ A set of $70000$ single-channel images of handwritten digits (split into proportion of $6:1$ of train to test data). Their size is fixed to $28$x$28$ pixels (although, for the purpose of some experiments, depending on the use case, its size was changed) and content centered. Due to its simplicity (relatively small size of images, uncomplicated content), large sample size (around $7000$ instances per class) and portraying a real-world problem, it has become a default dataset for testing models.
    \item CIFAR-10 \cite{Krizhevsky2009}
    \\\\ A collection of $60000$ RGB\footnote{An additive color model that is operates on $3$ channels representing colors — red, green and blue.} photos ($50000$ in train and $10000$ in test sets), which are a curated subset of the \textit{80 Million Tiny Images}\footnote{The \textit{80 Million Tiny Images} dataset has been formally withdrawn on 29th June 2020 as a consequence of offensive images and derogatory terms as categories found in it. \cite{Birhane2021}} data. Content of the $32$x$32$ pixels images includes $10$ mutually exclusive classes of real-world objects, much more complex than in MNIST. Because of that, the downstream task is also more intricate and presents an assay for simple classical approaches.
    \item CIFAR-100 \cite{Krizhevsky2009}
    \\\\ CIFAR-100 is just an extended version of the CIFAR-10 dataset (with the same base data, image parameters and train-to-test split ratio) — it yields $100$ classes, but with only $600$ examples per category. As a result of that, only the most modern and convoluted architectures achieve an accuracy score of over $90\%$. It posed the greatest challenge both to representation learning and evaluation methods.
    % \item Custom datasets
    % \\\\ Sometimes, observational experiments require datasets with particular characteristics, such as size, dimensionality, distribution, etc. In an effort to present the concepts in a more coherent, concise or comprehensible way, additional datasets were used. The use of said datasets was always preceded with an adequate description.
\end{itemize}

\section{Observational study}
A diligent research expects an appropriate context and perspective, which is impossible to achieve without understanding how the used tools work in practice. For this reason, an observational study was conducted. In contrast to the \nameref{sec:quantitative}, instead of comparing representations, it focuses more on providing information on the practical use, as well as detecting possible benefits and risks of the methods.

\vspace{\baselineskip}
The implicit viewpoint of this work requires analyzing Bayesian representation learning methods — this means the Bayesian elements were treated with priority. Wherever possible, a reference to the feature learning was mentioned.

\vspace{\baselineskip}
Because all three used methods differ significantly on the 
different datasets and procedures, the study of each particular is unique and differs from one another. While the examples (especially visualizations) shown are a result of a single run of a given experiment, they are an approximation representing numerous runs executed during the preparation of this study.

\vspace{\baselineskip}
For the study brevity and clarity, only the MNIST dataset was used. If not specified otherwise, a stratified subset of $10000$, $28$x$28$ images served as a training set, with additional $1000$ test samples in a test set where needed.

\vspace{\baselineskip}
The results of this study can be found in \nameref{chapter:observational}.


\section{Quantitative study}\label{sec:quantitative}
In order to assess the quality of the representations on a downstream task, a quantitative study was conducted. It aims to answer a question about how good the representations truly are in a real-world application.

\vspace{\baselineskip}
All three pairs of methods were analyzed. At first, the representations were trained using specific methods. The training duration was $50$ epochs for deep models (neural networks and autoencoders) and $5000$ iterations for PCAs. In order to provide reliable results, training of each set of hyperparameters of every model was repeated $30$ times, with different random seed values.

\vspace{\baselineskip}
As the second step, all the representations were evaluated with 5 simple classifiers, each trained with 5-fold cross-validation: Logistic Regression, Support Vector Machine, Naive Bayes, Multilayer Perceptron and Decision Tree. The used metric was a \textit{F1-score}, which was micro-averaged over classes.

\vspace{\baselineskip}
In the end, all scores were gathered, mean and standard deviation were calculated and rounded to three decimal places. After that, the results were presented in \nameref{chapter:quantitative}. In each table, there are specified dataset, model and possible additional parameters, along with the F1-score values.

\vspace{\baselineskip}
In order to keep a balance between data used for learning representations and their evaluation, each test set of the datasets described in \autoref{sec:datasets} was designated to the latter task. That being said, the train parts became a dataset on their own. This means that in the feature learning stage they were split again to train, test, or validation sets according to the needs.

% 	\cleardoublepage
	
	\chapter{Observational study}\label{chapter:observational}
\thispagestyle{chapterBeginStyle}

\include*{observational/bpca}
\include*{observational/bnn}
\include*{observational/vae}

% 	\cleardoublepage
	
	\chapter{Quantitative study results}\label{chapter:quantitative}
\thispagestyle{chapterBeginStyle}
To make this chapter more legible and easier to understand, a few assumptions were:
\begin{itemize}
    \item Results shown in the tables are always sorted from the most simple, to the most complex dataset. They are denoted with their names capitalized.
    \item Names of the representation learning models were denoted with their abbreviations. For Principal Component Analysis and Neural Networks, their Bayesian counterparts got a \textit{B} (for \textit{Bayesian}) prefix; for Autoencoders, a \textit{V} (for \textit{Variational}) prefix.
    \item Classification results were presented in last five columns and denoted as follows:
    \begin{itemize}
        \item Logistic Regression — LR,
        \item Support Vector Machine — SVM,
        \item Naive Bayes — NB,
        \item Multilayer Perceptron — MLP,
        \item Decision Tree — DT.
    \end{itemize}
    \item Classifiers' scores were placed in their respective columns and denoted as \textit{mean±std}.
    \item The best score for each dataset was written in bold.
\end{itemize}

\section{Raw data}
Learned representations should facilitate the training process, in relation to a raw dataset. For this reason, evaluations needed to be run on unprocessed data first. This set a reference point for further experiments. \autoref{tab:orig} shows results of this run. The results are surprising — models that are the best for one dataset might perform the worst in the other and vice versa. This, however, confirms a correct choice of the datasets and classifiers for evaluation.

\include*{quantitative/tables/orig}


\include*{quantitative/bpca}
\include*{quantitative/bnn}
\include*{quantitative/vae}

\section{Conclusions}
Tested representation learning approaches do prove to be useful for downstream tasks of other classificators. Bayesian methods, besides the advantages enumerated in \autoref{chapter:methods} and \autoref{chapter:observational}, often bring better results than their conventional counterparts. This is partly due to the way representations are produced, which has higher information retention. The latent space itself also has better properties due to natural regularization. It is important, however, to carefully adjust hyperparameters, especially those responsible for the Bayesian inference.
% 	\cleardoublepage
	
	\chapter{Conclusions and further work}\label{chapter:conclusion}
\thispagestyle{chapterBeginStyle}
All the goals defined in \autoref{sec:scope} were met across this work, particularly, the main objective of introducing, analyzing and testing a cross-section of Bayesian representation learning models. In \autoref{chapter:background} an extensive literature review and theoretical approach to Bayesian representation learning were provided. \autoref{chapter:methods} presented theoretical background for the 3 chosen Bayesian representation learning methods, their observational analysis was conducted in \autoref{chapter:observational} and the quantitative study of the downstream task testing in \autoref{chapter:quantitative}.

\vspace{\baselineskip}
From the conducted experiments, important conclusions need to be made:

\begin{itemize}
    \item Bayesian treatment enables real quantification of uncertainty, which is unavailable for conventional methods. It encourages better model interpretability.
    \item Bayesian representation learning methods preserve more important information about the input. By naturally regularizing the latent space, they obtain its better properties. This directly translates to better results of representation evaluation on downstream tasks.
    \item Seemingly simpler Bayesian algorithms, such as BPCA, are not necessarily quicker. Often they do not provide a way for batch processing, which requires operating on a whole dataset at once, which is not optimal.    
    \item Some Bayesian models, similarly to conventional, might be sensitive to hyperparameters. If possible, an optimization is recommended for satisfactory results.
    \item Bayesian approach, while saving computational effort in a global scale (by learning distributions, and not point estimates), is more expensive from the perspective of training individual models. For analytical purposes, where large numbers of experiments need to be run, it is recommended to use well-established implementations, ideally GPU-enabled or at least concurrency-enabled.
\end{itemize}

\vspace{\baselineskip}
While diligent, the conducted study is not exhaustive. The topic would benefit from the following future work and improvements:
\begin{itemize}
    \item adding more datasets with different downstream tasks (e.g., regression),
    \item considering different aspects of measuring representations quality (e.g., disentanglement, geometric priors, etc.),
    \item testing more diverse and advanced models (semi-supervised, contrastive, etc.),
    \item implementing ablation study (both on model elements and data),
    \item further hyperparameter optimization, with an emphasis on prior distributions, but also improving VAE's quality.
\end{itemize}

% 0.2 of https://arxiv.org/pdf/2002.08791.pdf
	\cleardoublepage
	
	
	%%%%%%%%%%%%%%%%%%%%%%%%%%%%%%%%%%%%%%%%%%%%%%%%%%%%%%%%%%%%%%%%%%%%%%%%%%%%%%
	%%%%%%%%%%%%%%%%%%%%%%%%%%%%%%% BIBLIOGRAFIA %%%%%%%%%%%%%%%%%%%%%%%%%%%%%%%%%
	%%%%%%%%%%%%%%%%%%%%%%%%%%%%%%%%%%%%%%%%%%%%%%%%%%%%%%%%%%%%%%%%%%%%%%%%%%%%%%

	\pagestyle{bibliographyStyle}
	\bibliographystyle{ieeetr}
	\bibliography{literature}
	\nocite{*}
	\thispagestyle{chapterBeginStyle}
        \addcontentsline{toc}{chapter}{Bibliography}

	\cleardoublepage
	
	%%%%%%%%%%%%%%%%%%%%%%%%%%%%%%%%%%%%%%%%%%%%%%%%%%%%%%%%%%%%%%%%%%%%%%%%%%%%%%
	%%%%%%%%%%%%%%%%%%%%%%%%%%%%%%%%% DODATKI %%%%%%%%%%%%%%%%%%%%%%%%%%%%%%%%%%%%
	%%%%%%%%%%%%%%%%%%%%%%%%%%%%%%%%%%%%%%%%%%%%%%%%%%%%%%%%%%%%%%%%%%%%%%%%%%%%%%
    
    % \listoflistings
    
% 	\appendix
% 	\pagestyle{appendixStyle}
	
% 	\input{appendices/appendix}
% 	\cleardoublepage

\end{document}

